% ===== §10 =====
\section{The Three-Register Dialectic of the Relational Real}
\label{sec:manifesto}

\noindent
The antithesis does not proceed by adding a seventh theory to the six. It proceeds by stating the framework within which the six are to be reread, each granted its truth and shown its limit, and within which the resonance the previous section identified can at last be activated. This section states that framework. It is a variant of the Lacanian theory of the three registers, and it is arrived at by four determinate revisions of that theory, each of which will be set out with its definition, its difference from orthodox Lacan, its justification, its positive claim, and what it yields for the theory of culture. The framework is named, at the close, the \emph{three-register dialectic of the relational Real}.

\subsection*{Why Lacan, and why revised}

The choice of Lacan as the host theory is not eclectic but motivated by a single property: integrative capacity. The Lacanian triad of the Imaginary, the Symbolic, and the Real is one of the few conceptual apparatuses that can hold, within a single structure, everything this paper must relate---the non-symbolizable ground of experience (the Real), the order of language, law, and meaning (the Symbolic), the field of image and identification (the Imaginary), the organization of power (through the phallic signifier within the Symbolic), the economy of drive and its lack (through the objet petit a and the constitutive lack), and the phenomenology of the subject's experience across all three \citep{lacan2006ecrits,lacan1977xi}. Each of the six traditions grasps one facet of this and absolutizes it: the interpretive tradition takes the Symbolic alone, the power tradition takes the Symbolic-as-power alone, the transmission tradition takes a sedimented Imaginary-and-Symbolic and freezes it. The virtue of the Lacanian triad is that it can house them all at once and show each as partial. That is the reason to adopt it. The reason to revise it is that, left orthodox, it is a theory of the individual subject and a static topology, and this paper needs a theory of culture, of relation, and of historical motion. The four revisions convert the one into the other.

\subsection*{Revision One: the Real has a neglected relational form}

\emph{Definition.} Orthodox Lacan gives the Real chiefly a negative and individual face: the Real is what resists symbolization absolutely, the traumatic remainder, the ineffable that the signifier cannot reach, encountered by the subject as the abyss around which anxiety turns \citep{lacan1992vii}. The first revision does not deny this face. It claims that the Real has, besides its negative traumatic form, a \emph{positive relational form}: the non-symbolizable coupling generated between two subjects, the ``us'' that no signifier exhausts, is Real in the precise sense that it exceeds symbolization, yet it is not a trauma and not an abyss but a bond. Both forms of the Real coexist; the revision recovers the second, which the tradition, fixed on the first, left in shadow.

\emph{Difference from orthodox Lacan.} The revision is deliberately moderate. It does not assert that all Reality is relational, which would be a wholesale ontological reconstruction the paper does not need and would have to defend against every instance of purely individual traumatic Real. It asserts only that alongside the individual traumatic Real there is a relational Real, a positive unsymbolizable borne in the coupling of subjects. This is consonant with a claim developed elsewhere by the present author, that relation is the unsymbolizable in its positive form.

\emph{Justification.} The interpretive tradition (\xref{sec:hermeneutic}) locks culture in the Symbolic, and the power tradition (\xref{sec:power}) locks it in power, precisely because neither has a relational Real to appeal to. Something in intimate culture is neither a meaning nor an encoding of force---the unsayable, unownable ``between'' of two people---and only a relational Real can name it. \emph{Claim.} The relational form of the Real is the non-symbolizable core of a relational coupling; it is the anchor of culture outside both sign and power, standing beside the traumatic Real rather than replacing it. \emph{Cultural yield.} Endogenous culture has, in the register of the Real, a root that power---whose organ, the phallic signifier, operates in the Symbolic---cannot reach.

\subsection*{Revision Two: the Real is not \emph{completely} symbolizable}

\emph{Definition.} The relational Real, like the traumatic Real, resists complete symbolization: it always leaves a remainder that no sign exhausts and no accounting captures. This revision is a deliberate \emph{retention} of what orthodox Lacan gets right, imported as a discipline against a naive constructivism that would make culture wholly a matter of signs. \emph{Difference from orthodox Lacan.} Almost none; the fidelity is the point.

\emph{Justification.} Were the Real completely symbolizable, culture could in principle be fully encoded, fully captured, fully commodified, and the resistance that both experience and justice testify to would be unintelligible. The remainder must be preserved. \emph{Claim.} The Real core of culture resists symbolization and therefore resists commodification: this is the ontological ground of what a later section calls the decoherence of capture---the symbolic surface of a bond can be forged and sold while its Real coherence is drawn off and lost, because the surface is symbolizable and the core is not. \emph{Cultural yield.} Intimate culture cannot be completely captured by capital; the claim now has a ground in the structure of the Real rather than resting on sentiment.

\subsection*{Revision Three: culture spans all three registers, none dispensable}

\emph{Definition.} Culture is not a Symbolic phenomenon alone. Every cultural phenomenon is constituted \emph{at once and necessarily} across the three registers: in the Imaginary (image, identification, the aesthetic face), in the Symbolic (language, rule, institution), and in the Real (the non-symbolizable relational core). A ``culture'' present in only one register is a mutilated abstraction; the presence in each register is substantial and non-derivative.

\emph{Difference from orthodox Lacan.} The triad, in Lacan, structures the \emph{subject}; here it structures \emph{culture}, a relational object, extending the device as the intimacy series has already extended it to relational trust across the three registers. \emph{Justification.} The transmission tradition seizes an Imaginary-Symbolic sediment, the interpretive tradition the Symbolic web, the power tradition the Symbolic-as-power; each absolutizes one register. To span all three is the sublation of their several one-sidednesses. \emph{Claim.} Culture has a substantial mode of being in each register; any theory that treats only one is partial. \emph{Cultural yield.} This directly answers the power tradition (\xref{sec:power}): if culture were only Symbolic, Foucault and Gramsci would win, for the phallic signifier organizes the Symbolic; but because culture has a substantial, non-derivative existence in the Real---where, by Revisions One and Two, a relational and non-symbolizable core resides---culture possesses an anchor the organ of power cannot reach, and is therefore not locked shut by power. This is the ground the sublation of power (\xref{sec:powersublation}) will stand on.

\subsection*{Revision Four: the three registers are not static but in materially driven, dialectical motion}

\emph{Definition.} In orthodox Lacan of the late topology, the three registers are knotted in a Borromean structure, a synchronic and relatively stable form in which each ring is necessary and none superior, the whole holding the subject together \citep{lacan1975rsi}. The fourth revision sets this structure into \emph{motion}, and the motion has a determinate form with three components. The first is the \emph{direct dialectic between the registers}: the Symbolic continually works to capture the Real---to name it, to rule it, to render the unsayable ``between'' into a sign, a rite, at the limit a commodity---while the Real continually overflows the Symbolic, exceeding every naming and returning as the remainder that requires a new form; the Imaginary is the mediating field and the battleground on which this contest is fought. The second is the \emph{deep dialectical form}, Hegelian rather than merely oppositional: the three registers may be read as moments of a self-unfolding of relational spirit---the Imaginary as immediacy (the unmediated image, the specular identification), the Symbolic as mediation and determinate negation (the image cut and negated by the sign), the Real as the negation of the negation (what returns, after the failure of symbolization, as a higher unsymbolizable that is now a mediated Real and not the original one). The third is the \emph{material motor}: this motion is not a self-movement of pure spirit, which would be idealism, but is driven by material conditions---which signifiers the Symbolic has to capture the Real with, and in what form the Real overflows, are conditioned by the forces and relations of production, so that media determine the available symbolic forms and material conditions determine the possible forms of relation. This is the joint at which historical materialism and Lacan are welded: the dialectic of the registers has a material engine.

\emph{Difference from orthodox Lacan.} The synchronic, relatively stable Borromean topology is converted into a diachronic, materially driven contradiction; the Real ceases to be a static resistant remainder and becomes an active moment that overflows and demands new symbolization.

\emph{Justification.} A static triad yields the \emph{structure} of culture but not its \emph{history}. Only registers in materially driven contradiction, with the Real perpetually overflowing, can ground---ontologically---the claims that culture has no final truth, is never static, and develops.

\emph{Programmatic claim} (a disciplined historical dialectic, not the slogan that culture has ``no truth''): {\itshape The truth of culture is historical and non-terminal. Every cultural form is, under specific material conditions, one capture by the Symbolic of the relational Real, and is to be evaluated dialectically: one acknowledges its progressive significance under its own historical conditions, and one recognizes the limit it discloses when conditions change and the Real overflows anew. Culture therefore develops---not by converging on an absolute truth, but by the historical sublation of its own prior forms, through the materially driven motion of the three registers. To criticize a cultural form and to acknowledge its former progressiveness are two faces of one dialectical evaluation.}

\emph{Cultural yield.} Cultural evolution is the diachronic unfolding of this contradiction; every form is an incomplete capture of the relational Real; capital's commodification is the extreme form of Symbolic capture and necessarily meets the Real's resistance; the encoding of power likewise cannot lock culture shut. This revision supplies the ontological ground for the dialectical sublation of power (\xref{sec:powersublation}) and the dialectic of cultural evolution (\xref{sec:evolution}). By way of illustration: in criticizing an old culture---a rite, a hierarchy of rank---one nonetheless acknowledges, dialectically, its progressive meaning under the conditions that produced it.

\subsection*{The position, and the chain it unlocks}

Taken together the four revisions constitute a nameable position, the \emph{three-register dialectic of the relational Real}: a Lacanian variant in which the Real is granted a relational form, in which that Real remains not completely symbolizable, in which culture spans all three registers without any being dispensable, and in which the registers move in a materially driven dialectic. The position unlocks a single chain that is the theoretical spine of the whole paper: the dialectic of the three registers means the Real perpetually overflows the Symbolic; because the Real overflows, culture cannot be locked shut by either power or capital, both of which operate by Symbolic capture; because it cannot be locked shut, culture can develop historically and dialectically; and because its development is bound at once by the resistance of the Real and by the criterion of justice, it has direction without a terminus---no final truth, but not mere drift either. Along this chain the psychoanalytic line is promoted from a mere ``ground in drive'' to the very skeleton of the framework, the three registers, within which drive, power, and the Real each find their place; and along this same chain the lines of historical materialism and of the Hegelian dialectic are welded to that skeleton, the one as the material motor of the registers' motion, the other as its form. With the framework stated, the antithesis can do its proper work: expose the presupposition the six traditions share, and invert it.
